\paragraph{The multiplier for the trace ideal and its top coefficient}
Write \(S_{\mathbb T}=\mathbb T[(X_j)_{j\in J}]\), viewed as a
subring of the product of the polynomial rings for the systems \(x\).
Write \(N\) for the family of representations \(N_x\), and extend
coefficients to \(\mathcal A_{\mathbb T}=S_{\mathbb T}[G_E]\).
The set
\[
 \mathscr I=\{u\in\mathcal A_{\mathbb T}:
       \operatorname{tr}N(uv)\in\mathfrak a
       \text{ for all }v\in\mathcal A_{\mathbb T}\}
\]
is a two-sided ideal. Right multiplication is absorbed into \(v\),
and left multiplication uses cyclicity of trace. It contains the
kernel of \(N\). For a tied row \(b\), the contributions to
\(b\cdot\mathbf y\) from the common kernel, \(p_\ast\mathcal I_\mathrm{aug}\),
and the inertia--Frobenius relations belong to \(\mathscr I\).
The remaining term is \(m\rho_0\) with \(m\) in the augmentation ideal.
Express \(m\) as a linear combination of \(g-1\) and replace
each \(g-1\) by \(\sum_i z_i(g)(g_i-1)\).
The change is in the kernel of \(N\), so \(m\) can be taken in
the left ideal generated by the \(g_i-1\).
Thus for each square matrix \(B\) of tied rows, some matrix \(B_1\)
satisfies
\[
 B\mathbf y\equiv (B_1\mathbf y)\rho_0\pmod{\mathscr I}.
\]
\leanlinks{\href{\upstream/grh/src/New/Congruences/FirstOrder/RibetLemma.lean\#L105}{The decomposition of tied rows modulo the trace ideal}
  \leanlinksep \href{\upstream/grh/src/New/Congruences/FirstOrder/RibetMethod.lean\#L360}{The two-sided trace ideal}}

We use that characteristic polynomials of block matrices have
coefficients in \(S_{\mathbb T}\). Amitsur's theorem says that if
the elements with characteristic polynomial over a coefficient
subring \(R\) span an algebra over \(R\), then all elements of
the algebra have that property.
The proof works in the ring with \(T^e=0\) for arbitrary \(e\).
If \(\gamma\) is an \(R[T]\)-multiple of a spanning element, then
\[
 (1-T\beta)(1-T^{j+1}\gamma)
   =1-T\bigl(\beta+T^j\gamma-T^{j+1}\beta\gamma\bigr).
\]
The determinant on the left is in \(R[T]/(T^e)\) provided
\(\det(1-T\beta)\) is there. If additions by \(T^{j+1}\) times
arbitrary elements are already allowed, the last correction on
the right can be removed. Descend from \(j=e\), sum over spanning
elements, and take \(j=0,\beta=0\). Varying \(e\) gives every
coefficient of the reversed characteristic polynomial.
Here the group elements span the group algebra and have their
characteristic coefficients in \(S_{\mathbb T}\).
For block matrices, the matrices with a single nonzero block row
span. Such a matrix is a block column times a block row.
The product in the reverse order is the original diagonal block.
Sylvester's determinant identity gives the two characteristic
polynomials up to powers of the variable, so the same descent
applies to block matrices. A further coefficient variable \(Y\)
is allowed by applying the argument over \(R[Y]\), where constant
matrices span the matrices over the polynomial ring.
\leanlinks{\href{\upstream/grh/src/Lib/General/LinearAlgebra/Eigenvalues.lean\#L647}{The descent on characteristic polynomial coefficients}
  \leanlinksep \href{\upstream/grh/src/New/Congruences/FirstOrder/RibetMethod.lean\#L44}{Application to block matrices}}

Apply this to \(B-YB_1\), and write \(P_B(Y)\in S_{\mathbb T}[Y]\)
for its block determinant through \(N\).
Apply Cayley--Hamilton to the characteristic polynomial
\(Q(Z)=ZQ_1(Z)+P_B(Y)\) of \(-(B-YB_1)\).
It gives a group-ring matrix \(B'=Q_1(-(B-YB_1))\) for which
\(B'(B-YB_1)=P_B(Y)I\) holds coefficientwise through \(N\).
Read the equality in \(\mathcal A_{\mathbb T}/\mathscr I\).
Let \(Y\) act on this quotient by right multiplication by \(\rho_0\),
and the group-ring coefficients by left multiplication.
The two actions commute by associativity, so the row congruence
says that \(B-YB_1\) in this action kills \(\mathbf y\).
Hence \(y_iP_B(\rho_0)\in\mathscr I\). For every \(w\) in the
left ideal generated by \(y_i=g_i-1\), this gives
\(wP_B(\rho_0)\in\mathscr I\).
\leanlinks{\href{\upstream/grh/src/New/Congruences/FirstOrder/RibetMethod.lean\#L150}{The adjugate and the variable acting by right multiplication}}

Set \(f_B=\operatorname{tr}N(Ku_0P_B(\rho_0))\), where \(u_0\)
is the group-ring element for the character indexed by \(a=0\).
The trace belongs to \(S_{\mathbb T}\) as a characteristic
polynomial coefficient. By Cayley--Hamilton again, some left
multiple of \(w\) acts through \(N\) as the scalar matrix
\(\det N(w)\). Multiply that left multiple by \(P_B(\rho_0)\)
to get an element of \(\mathscr I\). Multiply on the left by
\(Ku_0\) and take traces to get \(f_B\det N(w)\in\mathfrak a\).
In particular \(w=-\sum_i z_i(g)(g_i-1)\) has the same image
under \(N\) as \(1-g\), so \(f_Bd_g\in\mathfrak a\).

On \(x\in\mathcal X\), \(Ku_0\) acts by \(C\) on the trivial
character and by zero elsewhere. Since \(\rho_0\) vanishes on
that character, the trace is \(C P_B(0)(x)\).
The matrix
\[
 B=(1-K)B_0+KB_{\mathbf i}
\]
consists of tied rows and reads as \(B_0\) on the main block
and \(B_{\mathbf i}\) on the characters. For this \(B\),
\[
 f_B(x)=C D_x(X)c_0^{t'}
                \prod_{a=0}^{t'-1}\det B_{a,i(a)}.
\]
Sum with weights \(\prod_a w_{a,i(a)}\) to form \(\mathcal D\).
Then \(d_g\mathcal D\in\mathfrak a\), while on \(\mathcal X\)
we obtain \(\mathcal D(x)=C c_0^{t'}D_x(X)\prod_a d_a(X)\).
\leanlinks{\href{\upstream/grh/src/New/Congruences/FirstOrder/RibetMethod.lean\#L470}{The determinant multiplier from the trace}
  \leanlinksep \href{\upstream/grh/src/New/Congruences/FirstOrder/RibetLemma.lean\#L218}{Combining the main-block and character rows}}

In \(\prod_a d_a\), replace \(d_a\) by \(d_a+\xi_a\) in increasing
order of \(a\). The term added at step \(a\) is a multiple of
\(X_{<a}\xi_a\) and hence lies in \(\mathfrak n\).
Therefore \(\prod_a d_a\equiv X_J\prod_a\Gamma_a\bmod\mathfrak n\).
The coefficients at the proper divisors of \(X_J\) vanish.
The coefficient of \(X_J\) is
\[
 (C c_0^{t'}D_x\prod_a\Gamma_a)(0).
\]
At \(\pi\) its valuation is at most
\[
 b_{\mathrm{main}}+\sum_a b_a+
 t'\!\left(\sum_v v_\ell(o_v(1))+2\#s\,v_\ell(C)\right)
       +(\#J+1)v_\ell(C).
\]
The two contributions involving \(\#s\) come from \(c_0^{t'}\)
and the individual terms in \(\prod_a\Gamma_a\).
\leanlinks{\href{\upstream/grh/src/New/Congruences/FirstOrder/FittingBlock.lean\#L531}{Coefficients after combining the corrections}
  \leanlinksep \href{\upstream/grh/src/New/Congruences/FirstOrder/RibetLemma.lean\#L651}{The valuation of the top coefficient}}
